\section{Experimental Configuration Details}
\label{app:experimental-config}

Table~\ref{tab:dataset-spec} describes the generation grid and the two task families. The archived field \texttt{n\_if} appears as \texttt{if} in the diagnostic exports, while the paper defines the families through \mbox{\texttt{n\_if}$=0$} and \mbox{\texttt{n\_if}$>0$}. All archived configurations record generator mode \texttt{d}, including zero-condition instances, but this historical field does not change the analysis partition. Table~\ref{tab:models} gives the serving configuration recovered from the logs. Since model and tokenizer revisions are unset, the checkpoint names do not establish an immutable weight revision.

\begin{table*}[t]
\centering
\caption{Generation grid and task-family partition used in the reported experiment.}
\label{tab:dataset-spec}
\scriptsize
\begin{tabularx}{\textwidth}{p{0.22\textwidth}p{0.24\textwidth}X}
\toprule
Category & Parameter & Value \\
\midrule
Dataset size & Java classes & 1,232 classes \\
& Questions per model & 46,656: 11,664 unguarded and 34,992 condition-aware \\
& Question polarity & 23,328 reachable and 23,328 unreachable \\
\midrule
Graph generation & \path{class_size} & $\{100,250,500\}$ methods \\
& \path{depths} & $\{1,4,7,10,13,16,19,22,25\}$ \\
& Local tree structure & One structured execution tree spanning all generated methods; branching probability $0.65$ \\
& \path{useful_var_spawn_probability} & 0.2 \\
\midrule
Generated code & \path{n_comment} & $\{0,1\}$ comment lines per method \\
& \path{n_vars}, \path{n_loops}, \path{n_params} & $\{0,1,2\}$ \\
& \path{n_if} & $\{0,1,3,5\}$ \\
\midrule
Factorial design & Code configurations & $3\times2\times3\times3\times4\times3=648$ \\
& \path{n_questions} & Four reachable and four unreachable per distance and configuration \\
\midrule
Task semantics & Unguarded (\mbox{\texttt{n\_if}$=0$}) & No generated condition guards a call \\
& Condition-aware (\mbox{\texttt{n\_if}$>0$}) & Only executable calls, including calls whose guards evaluate to true, contribute to reachability \\
\bottomrule
\end{tabularx}
\end{table*}

\begin{table*}[t]
\centering
\caption{Model and serving configurations used for both task families. BF16 denotes the activation dtype, while FP8 denotes the quantization. The logs do not pin checkpoint or tokenizer revision hashes.}
\label{tab:models}
\scriptsize
\begin{tabularx}{\textwidth}{p{0.26\textwidth}p{0.08\textwidth}p{0.13\textwidth}X}
\toprule
Checkpoint & Groups & Budget & Serving configuration \\
\midrule
\texttt{google/gemma-4-26B-A4B-it} & U / C & 16,384 & vLLM 0.24.0; BF16/FP8; 262,144-token context; TP=1; \texttt{gemma4} reasoning parser \\
\texttt{Qwen/Qwen3.6-35B-A3B} & U / C & 16,384 & vLLM 0.24.0; BF16/FP8; 262,144-token context; TP=1; \texttt{qwen3} reasoning parser \\
\bottomrule
\end{tabularx}
\end{table*}


\section{Diagnostic Definitions and Populations}
\label{app:diagnostic-definitions}

Let $R(s)$ denote the set of methods executable from $s$, and $T$ the returned sequence. A guard is considered exposed when its containing method appears in $T$. We attribute a false-guard error when a newly introduced method has no executable predecessor among earlier trace methods and exactly one earlier guard explains its introduction. When several guards could explain it, attribution is undefined. This allows us to identify unsupported inclusions through false guards, without observing all possible condition-evaluation errors, such as ignoring a true branch.

\vspace{0.2cm}

T01 divides the number of attributed guard errors by the number of exposed call-controlling guards, with both counts summed over eligible responses. A gives the fraction of attributed errors whose guarded callees can reach the target through executable edges, while B gives the fraction of eligible responses containing at least one such error. For negative questions, T04 measures the mean fraction of $R(s)$ omitted, and T05 measures the incidence of traces containing an existing method outside $R(s)$. These defects can occur in the same trace, so the measurements do not form separate categories.

\vspace{0.2cm}

For structurally admissible positive traces, T06 measures target omission, while T07 measures continuation after the target among traces containing it. Structural admissibility requires a source-first sequence with no duplicate or unsupported method introduction. T08 measures unknown method names, T09 gives the relative position of the first structural error in nonempty defective traces, and T11 measures duplicates in eligible nonempty traces. T13 measures mean source-anchored coverage, corresponding to the fraction of $R(s)$ linked back to the source through executable introductions. This does not represent positive-query success, since a valid positive trace can stop at the target without covering the whole region.

\vspace{0.2cm}

For a guard $g$ whose variable origin is resolved, $h(g)$ counts the method boundaries back to the local declaration supplying its tested value. D1 averages $h(g)$ over exposed guards, while D3 restricts this average to guards with an attributed false-guard error and D5 to the remaining exposed guards in unambiguous traces. Literal conditions, unresolved origins, and ambiguous lineages are excluded. D1$^*$ uses the attribution population shared by D3/D5, allowing its numerator and denominator to be partitioned into those of D3 and D5. D9 averages $h(g)$ over guards in each reference depth-first exploration before averaging across questions. This reference traversal is a diagnostic convention and does not impose an exploration order on the model.

\vspace{0.2cm}

Table~\ref{tab:difficulty-populations} describes these different weights and exclusions. D1 differs from D1$^*$ because it also includes resolved exposures in traces with ambiguous error attribution. Table~\ref{tab:diagnostic-populations} gives the populations used for the other diagnostics, with missing measurements excluded from interpretations of success. The hashes accompanying the reproducible aggregates identify the source export and its metric catalogue.

The main text focuses on omissions, inclusions outside the executable region, and value-origin distance. Table~\ref{tab:error-taxonomy} retains the broader response checks, whose incidences overlap and do not partition failures. Table~\ref{tab:trace-diagnostics} separates outcomes by polarity, while Table~\ref{tab:value-origin} compares origin distances. Reference difficulty increases from 4.25 to 5.29 boundaries between \path{n_params}=1 and \path{n_params}=2. However, it is weighted by question, while the error comparison is weighted by guard occurrence.

\begin{table*}[!t]
\centering
\caption{Taxonomy used to classify trace errors. Counts will be added after the final experiment and dashboard analysis.}
\label{tab:error-taxonomy}
\tiny
\setlength{\tabcolsep}{2pt}
\begin{tabularx}{\textwidth}{p{0.15\textwidth}p{0.27\textwidth}p{0.21\textwidth}p{0.055\textwidth}p{0.055\textwidth}p{0.15\textwidth}}
\toprule
Error family & Definition & Detection method & $n$ & \% & Example \\
\midrule
Syntactic-call bias & The trace follows a visible call that cannot execute under the generated condition. & Compare predicted edges with the executable graph. & \TODO{n} & \TODO{\%} & false guarded call \\
Missing reachable nodes & A negative trace omits part of the source-reachable closure. & Compare predicted and reachable node sets. & \TODO{n} & \TODO{\%} & source only \\
Early stopping & A positive trace ends before the target. & Check the final predicted node. & \TODO{n} & \TODO{\%} & source only \\
Over-exploration & A positive trace continues after the target. & Check the final predicted node. & \TODO{n} & \TODO{\%} & past target \\
Invalid introduction & A method has no executable incoming edge from any earlier predicted node. & Check executable edges from the presented set. & \TODO{n} & \TODO{\%} & cross branch \\
Target hallucination & The trace contains a target or method name that is absent from the generated class. & Compare names with the method list. & \TODO{n} & \TODO{\%} & unknown name \\
Negative over-exploration & A negative trace explores methods outside the region reachable from the source. & Check the source-reachable region. & \TODO{n} & \TODO{\%} & sibling subtree \\
Wrong source & The first method in the trace is not the source specified in the question. & Check the first method. & \TODO{n} & \TODO{\%} & target first \\
Repeated methods & The trace repeats methods even though the call graph is acyclic. & Check for duplicates. & \TODO{n} & \TODO{\%} & $A,B,A$ \\
Invalid JSON / schema failure & The response cannot be parsed using the required output schema. & Check the parser status. & \TODO{n} & \TODO{\%} & missing field \\
\bottomrule
\end{tabularx}
\end{table*}

% Generated by python3 scripts/render_results.py; source: data/results/aggregates.csv.
\begin{table}[t]
\centering
\caption{Condition-aware scores by polarity (\%). Each polarity has 17,496 questions per model. T04 gives the mean omitted fraction, while T05 gives the incidence of traces containing an existing method outside the executable region. Both are negative-only diagnostics.}
\label{tab:trace-diagnostics}
\scriptsize
\setlength{\tabcolsep}{3pt}
\begin{tabular}{llrrrr}
\toprule
Model & Polarity & Answer & Trace & T04 & T05 \\
\midrule
Gemma & Positive & 82.42 & 69.27 & -- & -- \\
Gemma & Negative & 78.10 & 24.99 & 39.37 & 51.38 \\
Qwen & Positive & 92.64 & 81.57 & -- & -- \\
Qwen & Negative & 75.85 & 34.04 & 38.01 & 48.50 \\
\bottomrule
\end{tabular}

\end{table}

% Generated by python3 scripts/render_results.py; source: data/results/aggregates.csv.
\begin{table}[t]
\centering
\caption{Value-origin distance in method boundaries. D1$^*$ restricts D1 to the same unambiguous attribution population as D3/D5. D3 concerns attributed errors, while D5 concerns exposed guards without an attributed error. D9 is question-weighted reference difficulty. Counts appear in Table~\ref{tab:difficulty-populations}.}
\label{tab:value-origin}
\scriptsize
\setlength{\tabcolsep}{3pt}
\begin{tabular}{lrrrrr}
\toprule
Model & \path{n_params} & D1$^*$ & D3 & D5 & D9 \\
\midrule
Gemma & 1 & 3.72 & 5.41 & 3.59 & 4.25 \\
Gemma & 2 & 4.65 & 6.22 & 4.47 & 5.29 \\
Qwen & 1 & 3.72 & 5.15 & 3.58 & 4.25 \\
Qwen & 2 & 4.63 & 5.90 & 4.47 & 5.29 \\
\bottomrule
\end{tabular}

\end{table}

% Generated by python3 scripts/render_results.py; source: data/results/aggregates.csv.
\begin{table*}[t]
\centering
\caption{Difficulty populations within condition-aware questions. Each \texttt{n\_params} level has 11,664 questions. The denominator counts resolved guard occurrences for D1/D1$^*$/D3/D5, and eligible questions for D9. Exclusions include zero resolved occurrences, missing origins, or ambiguous attribution as applicable.}
\label{tab:difficulty-populations}
\scriptsize
\setlength{\tabcolsep}{3pt}
\begin{tabular}{llrrrrr}
\toprule
Model & \path{n_params} & Metric & Mean & Eligible questions & Denominator & Excluded \\
\midrule
Gemma & 1 & D1 & 3.984 & 11,149 & 161,390 & 515 \\
Gemma & 1 & D1$^*$ & 3.724 & 9,337 & 95,113 & 2,327 \\
Gemma & 1 & D3 & 5.411 & 2,784 & 7,057 & 8,880 \\
Gemma & 1 & D5 & 3.588 & 9,328 & 88,056 & 2,336 \\
Gemma & 1 & D9 & 4.251 & 11,592 & 11,592 & 72 \\
Gemma & 2 & D1 & 4.922 & 11,073 & 163,650 & 591 \\
Gemma & 2 & D1$^*$ & 4.653 & 8,837 & 83,662 & 2,827 \\
Gemma & 2 & D3 & 6.223 & 3,202 & 8,701 & 8,462 \\
Gemma & 2 & D5 & 4.471 & 8,821 & 74,961 & 2,843 \\
Gemma & 2 & D9 & 5.285 & 11,575 & 11,575 & 89 \\
\midrule
Qwen & 1 & D1 & 3.928 & 11,353 & 143,319 & 311 \\
Qwen & 1 & D1$^*$ & 3.720 & 9,716 & 91,240 & 1,948 \\
Qwen & 1 & D3 & 5.145 & 3,279 & 8,222 & 8,385 \\
Qwen & 1 & D5 & 3.579 & 9,707 & 83,018 & 1,957 \\
Qwen & 1 & D9 & 4.251 & 11,592 & 11,592 & 72 \\
Qwen & 2 & D1 & 4.789 & 11,342 & 145,326 & 322 \\
Qwen & 2 & D1$^*$ & 4.625 & 9,400 & 84,280 & 2,264 \\
Qwen & 2 & D3 & 5.898 & 3,604 & 9,366 & 8,060 \\
Qwen & 2 & D5 & 4.466 & 9,382 & 74,914 & 2,282 \\
Qwen & 2 & D9 & 5.285 & 11,575 & 11,575 & 89 \\
\bottomrule
\end{tabular}

\end{table*}

% Generated by python3 scripts/render_results.py; source: data/results/aggregates.csv.
\begin{table*}[t]
\centering
\caption{Condition-aware diagnostic aggregates. Values are percentages, except for T09, which is a median with no sum denominator. Eligibility depends on the metric, so exclusions also include the other polarity for negative/positive-only diagnostics. For T04/T13, the numerator sums per-question fractions, not node counts. The two FP rows use unambiguously attributable negative records and condition on whether a target-directed guard error is present.}
\label{tab:diagnostic-populations}
\scriptsize
\setlength{\tabcolsep}{3pt}
\begin{tabular}{llrrrrr}
\toprule
Model & Metric & Value & Numerator & Denominator & Eligible & Excluded \\
\midrule
Gemma & T01 & 4.34 & 16,253.00 & 374,856 & 29,507 & 5,485 \\
Gemma & T04 & 39.37 & 6,888.62 & 17,496 & 17,496 & 17,496 \\
Gemma & T05 & 51.38 & 8,990.00 & 17,496 & 17,496 & 17,496 \\
Gemma & T06 & 8.48 & 1,135.00 & 13,392 & 13,392 & 21,600 \\
Gemma & T07 & 1.13 & 138.00 & 12,257 & 12,257 & 22,735 \\
Gemma & T08 & 0.24 & 85.00 & 34,992 & 34,992 & 0 \\
Gemma & T09 & 28.57 & -- & -- & 13,043 & 21,949 \\
Gemma & T11 & 5.71 & 1,943.00 & 34,014 & 34,014 & 978 \\
Gemma & T13 & 44.98 & 15,289.65 & 33,990 & 33,990 & 1,002 \\
Gemma & A & 19.81 & 3,220.00 & 16,253 & 30,784 & 4,208 \\
Gemma & B & 9.65 & 2,971.00 & 30,784 & 30,784 & 4,208 \\
Gemma & FP with target-directed error & 89.73 & 1,721.00 & 1,918 & 1,918 & 33,074 \\
Gemma & FP without target-directed error & 1.36 & 170.00 & 12,536 & 12,536 & 22,456 \\
\midrule
Qwen & T01 & 5.34 & 17,874.00 & 334,698 & 30,526 & 4,466 \\
Qwen & T04 & 38.01 & 6,650.81 & 17,496 & 17,496 & 17,496 \\
Qwen & T05 & 48.50 & 8,485.00 & 17,496 & 17,496 & 17,496 \\
Qwen & T06 & 4.12 & 614.00 & 14,893 & 14,893 & 20,099 \\
Qwen & T07 & 0.06 & 8.00 & 14,279 & 14,279 & 20,713 \\
Qwen & T08 & 0.59 & 207.00 & 34,992 & 34,992 & 0 \\
Qwen & T09 & 30.00 & -- & -- & 11,227 & 23,765 \\
Qwen & T11 & 1.29 & 445.00 & 34,534 & 34,534 & 458 \\
Qwen & T13 & 41.59 & 14,296.55 & 34,377 & 34,377 & 615 \\
Qwen & A & 23.30 & 4,165.00 & 17,874 & 31,291 & 3,701 \\
Qwen & B & 12.22 & 3,824.00 & 31,291 & 31,291 & 3,701 \\
Qwen & FP with target-directed error & 92.57 & 2,367.00 & 2,557 & 2,557 & 32,435 \\
Qwen & FP without target-directed error & 0.46 & 55.00 & 12,074 & 12,074 & 22,918 \\
\bottomrule
\end{tabular}

\end{table*}


\section{Detailed Result Curves}
\label{app:result-details}

Figure~\ref{fig:results-guards-parameters} shows the marginal profiles for \path{n_if} and \path{n_params}, while Table~\ref{tab:guard-mechanisms} distinguishes the rate of detected guard errors from their target-directed consequences. T01 decreases as \path{n_if} rises, although path scores deteriorate and A increases. This can occur because additional conditions increase the exposure denominator, while the returned traces determine which guards are exposed. A smaller T01 ratio therefore does not establish better condition evaluation, and counts must be interpreted together with eligibility.

\vspace{0.2cm}

Figure~\ref{fig:results-distance} separates signed distance levels, allowing the two polarities to be examined individually. Figure~\ref{fig:results-context} also shows that higher negative answer correctness can coexist with lower path correctness. These curves provide additional information about where the decision and its evidence diverge.

% Generated by python3 scripts/render_results.py; source: data/results/aggregates.csv.
\begin{figure*}[t]
\centering
\begin{tikzpicture}
\begin{groupplot}[rb score axis, group style={group size=2 by 1, horizontal sep=1.3cm}]
\nextgroupplot[title={(a) \texttt{n\_if}},xmin=-0.2,xmax=5.2,xtick={0,1,3,5},xlabel={Generated guards ($\mathtt{n\_if}$)},legend to name=fig-results-guards-parameters-legend,legend columns=4]
\addplot[rb gemma,rb answer] coordinates {(0,96.0477) (1,89.5319) (3,78.3779) (5,72.8738)};
\addlegendentry{Gemma: answer}
\addplot[rb gemma,rb path] coordinates {(0,82.2274) (1,59.3879) (3,44.5473) (5,37.4486)};
\addlegendentry{Gemma: trace}
\addplot[rb qwen,rb answer] coordinates {(0,99.1941) (1,92.5840) (3,82.3045) (5,77.8464)};
\addlegendentry{Qwen: answer}
\addplot[rb qwen,rb path] coordinates {(0,96.3306) (1,67.7726) (3,57.0302) (5,48.6025)};
\addlegendentry{Qwen: trace}
\nextgroupplot[title={(b) \texttt{n\_params}},xmin=-0.1,xmax=2.1,xtick={0,1,2},xlabel={Parameters per method ($\mathtt{n\_params}$)},ylabel={}]
\addplot[rb gemma,rb answer] coordinates {(0,91.2380) (1,75.7888) (2,73.7569)};
\addplot[rb gemma,rb path] coordinates {(0,66.6152) (1,39.1718) (2,35.5967)};
\addplot[rb qwen,rb answer] coordinates {(0,97.4108) (1,78.4379) (2,76.8861)};
\addplot[rb qwen,rb path] coordinates {(0,90.5350) (1,43.5185) (2,39.3519)};
\end{groupplot}
\node[anchor=north,inner sep=0pt] at ($(group c1r1.south)!0.5!(group c2r1.south)+(0,-1.05cm)$) {\pgfplotslegendfromname{fig-results-guards-parameters-legend}};
\end{tikzpicture}
\caption{Scores by \mbox{\texttt{n\_if}} over all questions (a), and by \texttt{n\_params} within condition-aware questions (b). Other factors are balanced. The zero-value level of \mbox{\texttt{n\_if}} is the unguarded group.}
\Description{Two panels with correctness on the same zero-to-one-hundred percent scale. Circles denote Gemma and squares Qwen; solid lines denote answer correctness and dashed lines trace correctness.}
\label{fig:results-guards-parameters}
\end{figure*}

% Generated by python3 scripts/render_results.py; source: data/results/aggregates.csv.
\begin{table*}[t]
\centering
\caption{Guard diagnostics by \mbox{\texttt{n\_if}} (\%). A gives the target-directed share of attributed errors, while B gives the incidence of at least one such error in an eligible response. T05 applies only to the 5,832 negative questions per condition level. T01 uses the exposed-guard denominator, while A/B exclude ambiguous attributions.}
\label{tab:guard-mechanisms}
\scriptsize
\setlength{\tabcolsep}{3pt}
\begin{tabular}{llrrrrrr}
\toprule
Model & \path{n_if} & Answer & Trace & T01 & A & B & T05 \\
\midrule
Gemma & 1 & 89.53 & 59.39 & 9.21 & 9.24 & 6.07 & 49.35 \\
Gemma & 3 & 78.38 & 44.55 & 3.81 & 25.91 & 11.92 & 52.52 \\
Gemma & 5 & 72.87 & 37.45 & 2.07 & 36.41 & 11.62 & 52.28 \\
Qwen & 1 & 92.58 & 67.77 & 10.88 & 11.19 & 7.34 & 42.37 \\
Qwen & 3 & 82.30 & 57.03 & 4.70 & 31.01 & 15.10 & 49.49 \\
Qwen & 5 & 77.85 & 48.60 & 2.94 & 36.63 & 15.06 & 53.64 \\
\bottomrule
\end{tabular}

\end{table*}

% Generated by python3 scripts/render_results.py; source: data/results/aggregates.csv.
\begin{table}[t]
\centering
\caption{Condition-aware context diagnostics (\%). Pos./neg. denote final-answer scores, T01 gives the attributed-error/exposure ratio, and A gives the target-directed share of attributed guard errors. Each polarity score has 5,832 questions per \texttt{class\_size} level.}
\label{tab:context-mechanisms}
\scriptsize
\setlength{\tabcolsep}{3pt}
\begin{tabular}{lrrrrr}
\toprule
Model & \path{class_size} & Pos. & Neg. & T01 & A \\
\midrule
Gemma & 100 & 96.04 & 69.24 & 4.06 & 32.36 \\
Gemma & 500 & 66.03 & 87.64 & 4.53 & 8.90 \\
Qwen & 100 & 97.67 & 67.25 & 5.04 & 35.27 \\
Qwen & 500 & 86.68 & 84.67 & 5.63 & 13.62 \\
\bottomrule
\end{tabular}

\end{table}

% Generated by python3 scripts/render_results.py; source: data/results/aggregates.csv.
\begin{figure*}[t]
\centering
\begin{tikzpicture}
\begin{groupplot}[rb score axis, group style={group size=2 by 1, horizontal sep=1.3cm}]
\nextgroupplot[title={(a) Final-answer correctness},xmin=-27,xmax=27,xtick={-25,-13,0,13,25},xlabel={Signed generation distance},legend to name=fig-results-distance-legend,legend columns=2]
\addplot[rb gemma,rb answer] coordinates {(-25,76.0802) (-22,76.5947) (-19,77.1605) (-16,77.4691) (-13,76.8519) (-10,77.8292) (-7,78.1379) (-4,79.0123) (-1,83.7449)};
\addlegendentry{Gemma}
\addplot[rb gemma,rb answer,forget plot] coordinates {(1,97.5309) (4,89.0432) (7,85.3909) (10,82.6132) (13,80.8642) (16,78.6008) (19,77.2119) (22,74.8971) (25,75.6687)};
\addplot[rb qwen,rb answer] coordinates {(-25,73.6111) (-22,75.6687) (-19,74.6399) (-16,74.2798) (-13,75.0000) (-10,74.1255) (-7,77.3148) (-4,76.1317) (-1,81.8930)};
\addlegendentry{Qwen}
\addplot[rb qwen,rb answer,forget plot] coordinates {(1,99.9486) (4,97.8909) (7,95.7305) (10,93.5185) (13,92.5412) (16,89.6091) (19,88.6317) (22,88.3745) (25,87.5000)};
\addplot[rbRule,densely dotted,no marks,forget plot] coordinates {(0,0) (0,100)};
\nextgroupplot[title={(b) Trace correctness},xmin=-27,xmax=27,xtick={-25,-13,0,13,25},xlabel={Signed generation distance},ylabel={}]
\addplot[rb gemma,rb path] coordinates {(-25,11.6255) (-22,14.1461) (-19,16.0494) (-16,19.5988) (-13,20.8333) (-10,23.9198) (-7,28.4465) (-4,36.2654) (-1,54.0123)};
\addplot[rb gemma,rb path,forget plot] coordinates {(1,95.6790) (4,81.9959) (7,72.9938) (10,69.2387) (13,65.5350) (16,64.0432) (19,59.9794) (22,57.3560) (25,56.5844)};
\addplot[rb qwen,rb path] coordinates {(-25,22.9424) (-22,25.6173) (-19,27.5720) (-16,30.0926) (-13,31.5844) (-10,34.2593) (-7,37.9115) (-4,40.7922) (-1,55.5556)};
\addplot[rb qwen,rb path,forget plot] coordinates {(1,99.8971) (4,91.4095) (7,86.6770) (10,81.0185) (13,81.0700) (16,75.9259) (19,75.6687) (22,71.9136) (25,70.5247)};
\addplot[rbRule,densely dotted,no marks,forget plot] coordinates {(0,0) (0,100)};
\end{groupplot}
\node[anchor=north,inner sep=0pt] at ($(group c1r1.south)!0.5!(group c2r1.south)+(0,-1.05cm)$) {\pgfplotslegendfromname{fig-results-distance-legend}};
\end{tikzpicture}
\caption{Condition-aware scores by signed generation distance. Negative distances denote unreachable targets, while positive distances denote reachable targets. Each point has 1,944 questions per model. The dotted divider at zero separates the two populations, with no zero-distance query plotted.}
\Description{Answer and trace correctness by signed distance. The two polarity-specific curves stop on either side of zero. Model identity uses the same colors and markers as the other result figures.}
\label{fig:results-distance}
\end{figure*}

% Generated by python3 scripts/render_results.py; source: data/results/aggregates.csv.
\begin{figure*}[t]
\centering
\begin{tikzpicture}
\begin{groupplot}[rb score axis, group style={group size=2 by 1, horizontal sep=1.3cm}]
\nextgroupplot[title={(a) Reachable targets},xmin=80,xmax=520,xtick={100,250,500},xlabel={\texttt{class\_size}},legend to name=fig-results-context-legend,legend columns=4]
\addplot[rb gemma,rb answer] coordinates {(100,96.0391) (250,85.2023) (500,66.0322)};
\addlegendentry{Gemma: answer}
\addplot[rb gemma,rb path] coordinates {(100,79.1324) (250,71.6049) (500,57.0645)};
\addlegendentry{Gemma: trace}
\addplot[rb qwen,rb answer] coordinates {(100,97.6680) (250,93.5700) (500,86.6770)};
\addlegendentry{Qwen: answer}
\addplot[rb qwen,rb path] coordinates {(100,83.8649) (250,82.5789) (500,78.2579)};
\addlegendentry{Qwen: trace}
\nextgroupplot[title={(b) Unreachable targets},xmin=80,xmax=520,xtick={100,250,500},xlabel={\texttt{class\_size}},ylabel={}]
\addplot[rb gemma,rb answer] coordinates {(100,69.2387) (250,77.4177) (500,87.6372)};
\addplot[rb gemma,rb path] coordinates {(100,32.2702) (250,25.3086) (500,17.3868)};
\addplot[rb qwen,rb answer] coordinates {(100,67.2497) (250,75.6344) (500,84.6708)};
\addplot[rb qwen,rb path] coordinates {(100,36.3512) (250,34.0364) (500,31.7215)};
\end{groupplot}
\node[anchor=north,inner sep=0pt] at ($(group c1r1.south)!0.5!(group c2r1.south)+(0,-1.05cm)$) {\pgfplotslegendfromname{fig-results-context-legend}};
\end{tikzpicture}
\caption{Condition-aware scores by class size, separated by polarity. Each point has 5,832 questions per model. Negative answer correctness rises while negative trace correctness falls.}
\Description{Correctness by class size in two polarity-specific panels. Solid answer lines and dashed trace lines use circles for Gemma and squares for Qwen, with a shared legend and percentage scale.}
\label{fig:results-context}
\end{figure*}


\vspace{0.2cm}

Table~\ref{tab:token-observability} reports persisted token counts, showing that condition-aware responses contain more reasoning tokens on average than unguarded responses despite lower primary scores. Mean reasoning usage on condition-aware questions is 45.72\% of the requested budget for Gemma and 38.43\% for Qwen. However, these means do not show whether individual failures exhausted their budget, and tokenizer differences prevent interpreting token ratios as equal compute. This experiment uses one requested budget and does not establish whether increasing it would resolve the observed failures.

% Generated by python3 scripts/render_results.py; source: data/results/aggregates.csv.
\begin{table*}[t]
\centering
\caption{Mean persisted reasoning and answer tokens, counted with each model's tokenizer. Usage gives reasoning tokens relative to the requested 16,384-token budget. These measurements describe output length rather than wall-clock compute or a budget intervention.}
\label{tab:token-observability}
\scriptsize
\setlength{\tabcolsep}{3pt}
\begin{tabular}{llrrr}
\toprule
Model & Group & Reasoning tokens & Usage (\%) & Answer tokens \\
\midrule
Gemma & Unguarded & 4,345 & 26.52 & 396 \\
Gemma & Condition-aware & 7,490 & 45.72 & 585 \\
Qwen & Unguarded & 3,388 & 20.68 & 83 \\
Qwen & Condition-aware & 6,296 & 38.43 & 96 \\
\bottomrule
\end{tabular}

\end{table*}


\section{Published Benchmark Comparisons}
\label{app:external-comparison}

Table~\ref{tab:external-comparison} gives scores reported in official model cards for the named checkpoints, checked on October~6, 2026. Both model columns come from Qwen's card, with Google's card also confirming Gemma's LiveCodeBench score~\cite{qwen2026qwen36,google2026gemma4}. These values therefore describe published external results rather than additional runs in our experiment.

\begin{table*}[t]
\centering
\caption{Published external scores (\%) for the named checkpoints. Scores can be compared within rows, while protocols differ across rows and from \benchmark.}
\label{tab:external-comparison}
\scriptsize
\begin{tabularx}{\textwidth}{lrrX}
\toprule
Benchmark & Gemma & Qwen & Provenance and protocol qualification \\
\midrule
LiveCodeBench v6 & 77.1 & 80.4 & Qwen card, with Gemma's value also confirmed by Google. \\
SWE-bench Verified & 17.4 & 73.4 & Qwen card, using an internal agent scaffold with shell and file-edit tools. \\
SWE-bench Multilingual & 17.3 & 67.2 & Same vendor-reported scaffold, not our reachability protocol. \\
SWE-bench Pro & 13.8 & 49.5 & Qwen card, with a refined benchmark set applied to its baselines. \\
Terminal-Bench 2.0 & 34.2 & 51.5 & Qwen card, using the Harbor/Terminus-2 harness and the mean of five runs. \\
\bottomrule
\end{tabularx}
\end{table*}

% Flush appendix floats and restore full-height columns before references.
\clearpage
\nolinenumbers
\twocolumn
